\section{Fazit und Ausblick}
\subsection{Celal:}
Im Verlauf des Semesters habe ich viele wertvolle Erfahrungen und Erkenntnisse im Modul „Vernetzte Systeme“ gesammelt. Besonders deutlich wurde für mich, wie essenziell das Verständnis von verteilten Systemen, Netzwerktechnologien und Tools wie Docker für die Entwicklung stabiler und skalierbarer Anwendungen ist. Durch die praktische Umsetzung konnte ich das theoretische Wissen schrittweise vertiefen und festigen.

Ein entscheidender Aspekt war die Zusammenarbeit im Team. Im Vergleich zum letzten Projekt verlief die Teamarbeit dieses Mal deutlich besser. Wir haben unsere Erfahrungen aus vorherigen Projekten genutzt, uns gegenseitig motiviert und dadurch eine produktivere sowie konstruktivere Arbeitsatmosphäre geschaffen. Dennoch gab es auch Herausforderungen, die uns zwangen, kreative Lösungsansätze zu entwickeln und Kompromisse einzugehen. Diese Erfahrungen haben mir die Bedeutung von Flexibilität und kontinuierlicher Kommunikation in der Teamarbeit noch bewusster gemacht.

Besonders herausfordernd waren die technischen Probleme, die immer wieder auftraten. Einige davon haben mich regelrecht verzweifeln lassen, doch letztendlich konnte ich sie mit viel Mühe lösen. Ein großes Hindernis war Redis, das zunächst nicht wie erwartet funktionierte. Ursprünglich hatten wir die Session-Verwaltung über MariaDB realisiert, sind dann aber auf Redis umgestiegen, weil wir Sticky Sessions implementieren wollten. Auch kleine Fehler haben uns oft stundenlang beschäftigt – winzige Unstimmigkeiten, wie ein falsches Zeichen an der falschen Stelle oder die richtige Konfiguration an der falschen Position, haben uns enorm aufgehalten. Doch durch mehrfaches Testen, Debuggen und systematisches Analysieren konnten wir diese Probleme schließlich überwinden. Dabei habe ich gelernt, mich kritisch mit Problemen auseinanderzusetzen, anstatt vorschnell eine Lösung zu erzwingen.

Ein weiteres großes Problem war der Lasttest, bei dem wir mit einer hohen CPU-Auslastung zu kämpfen hatten. Viele unserer ersten Ansätze erwiesen sich als ineffizient, sodass wir sie verwerfen und von vorne beginnen mussten. Durch diese Erfahrung habe ich erkannt, dass es manchmal sinnvoller ist, eine Lösung komplett neu aufzubauen, anstatt sich an einem fehlerhaften Konzept festzubeißen. Diese Denkweise hat mir geholfen, effizienter mit Herausforderungen umzugehen.

Durch die intensive Arbeit mit Docker habe ich großes Interesse an diesem Thema entwickelt. Mittlerweile beschäftige ich mich auch außerhalb dieses Projekts mit Docker, sehe mir YouTube-Tutorials dazu an und probiere eigene kleine Projekte aus. Ich finde es spannend, wie flexibel und leistungsfähig diese Technologie ist, und freue mich darauf, mein Wissen weiter zu vertiefen und in zukünftigen Projekten anzuwenden.

Insgesamt hat mir dieses Semester – trotz aller Höhen und Tiefen – ein tieferes Verständnis für vernetzte Systeme vermittelt. Die Kombination aus Teamarbeit, strukturiertem Vorgehen und bewusster Priorisierung meiner Aufgaben hat mich sowohl fachlich als auch organisatorisch weitergebracht. Diese Erkenntnisse werde ich in zukünftigen Projekten gezielt anwenden, um noch effizienter und zielgerichteter zu arbeiten.
